\section*{अध्याय \ref{ch:b1:atmospheric-circulation} --- वायुमंडल का सामान्य परिसंचरण}
% values: calc/ch28.py

\begin{solution}{exo:b1:atmospheric-circulation:1}
(क)~$f = 2\Omega\sin\varphi$: भूमध्य रेखा पर $0$, \ang{30} पर \qty{7.29e-5}{s^{-1}},
\ang{45} पर \qty{1.031e-4}{s^{-1}} और ध्रुव पर \qty{1.458e-4}{s^{-1}}।
(ख)~$fV = \num{1.031e-4} \times 10 = \qty{1.03e-3}{m.s^{-2}}$, $g$ का लगभग \num{1.05e-4}। (ग)~उत्तरी गोलार्ध में विक्षेपण गति के दाईं ओर है: उत्तर की ओर बहती पवन पूर्व की ओर धकेली जाती है।
\end{solution}

\begin{solution}{exo:b1:atmospheric-circulation:2}
(क)~$|\nabla p|/\rho = (400/\num{4e5})/1.225 = \qty{8.16e-4}{m.s^{-2}}$।
(ख)~$10/\num{8.16e-4} = \qty{12250}{s}$, \qty{3.4}{h}। (ग)~\ang{45} उ पर \qty{10}{m/s} की पवन का कोरिऑलिस त्वरण \qty{1.03e-3}{m.s^{-2}} है: दोनों एक ही आकार के हैं (अनुपात $0.79$)। चलना शुरू करने के कुछ घंटों के भीतर वायु पर एक कोरिऑलिस बल कार्य करने लगता है जो दाब-प्रवणता को संतुलित कर सकता है: पवन भू-विक्षेपी गति के निकट स्थिर हो जाती है, यहाँ $\num{8.16e-4}/\num{1.031e-4} = \qty{7.9}{m/s}$, प्रवणता के नीचे त्वरित होने के बजाय।
\end{solution}

\begin{solution}{exo:b1:atmospheric-circulation:3}
(क)~$f = 2 \times \num{7.292e-5} \times \sin\ang{40} = \qty{9.374e-5}{s^{-1}}$;
$|\vec v_g| = (400/\num{2.5e5})/(1.225 \times \num{9.374e-5}) = \qty{13.9}{m/s}$, या
$13.93/0.5144 = \qty{27.1}{kn}$। (ख)~उत्तरी गोलार्ध में पवन निम्न को अपनी बाईं ओर रखती है: निम्न दक्षिण में होने पर वह पश्चिम की ओर बहती है, पूर्व से (दिगंश \ang{90})। (ग)~स्टाफ़ स्टेशन-वृत्त से पूर्व की ओर इंगित करता है, पवन की ओर; \qty{27.1}{kn} \qty{25}{kn} तक पूर्णांकित होता है, दो पूर्ण पंख और एक आधा पंख, स्टाफ़ के दक्षिणावर्त पक्ष पर खींचे गए, यहाँ उसकी दक्षिणी ओर, निम्न की ओर।
\end{solution}

\begin{solution}{exo:b1:atmospheric-circulation:4}
(क)~\ang{15} उ: उत्तर-पूर्वी \omterm{def:b1:atmospheric-circulation:winds}{व्यापारिक पवनें}; \ang{15} द: दक्षिण-पूर्वी \omterm{def:b1:atmospheric-circulation:winds}{व्यापारिक पवनें}; \ang{45} उ और \ang{45} द: \omterm{def:b1:atmospheric-circulation:winds}{पछुआ पवनें}। (ख)~\ang{15} पर \omterm{def:b1:atmospheric-circulation:cells}{हैडली कोष्ठ} की सतही वायु भूमध्य रेखा की ओर लौटती है और कोरिऑलिस बल से पश्चिम की ओर मुड़ती है (उत्तर में दाईं ओर, दक्षिण में बाईं ओर): पूर्वी पवन, ऋणात्मक $u$। \ang{45} पर \omterm{def:b1:atmospheric-circulation:cells}{फ़ेरल कोष्ठ} की सतही वायु ध्रुव की ओर बढ़ती है और पूर्व की ओर मुड़ती है: पश्चिमी पवन, धनात्मक $u$। (ग)~चारों अक्षांशों में से \ang{30} उ पर दाब सबसे अधिक है, \qty{1016.8}{hPa}: \omterm{def:b1:atmospheric-circulation:cells}{उपोष्णकटिबंधीय उच्च} \ang{30} उ के निकट है (पूर्ण प्रोफ़ाइल उसका शिखर \ang{35} उ के निकट रखती है)।
\end{solution}

\begin{solution}{exo:b1:atmospheric-circulation:5}
(क)~$|\nabla p|/\rho = (200/\num{5e5})/1.225 = \qty{3.27e-4}{m.s^{-2}}$, जिसे
$f$ से भाग दें: \qty{2.6}{m/s} \ang{60} पर, \qty{2.9}{m/s} \ang{50} पर, \qty{4.5}{m/s}
\ang{30} पर, \qty{6.5}{m/s} \ang{20} पर, \qty{12.9}{m/s} \ang{10} पर और
\qty{25.7}{m/s} \ang{5} पर। (ख)~कोरिऑलिस बल एक पवन को समय $1/f$ की कोटि में मोड़ता है:
\qty{2.2}{h} \ang{60} पर, \qty{11}{h} \ang{10} पर, \qty{22}{h} \ang{5} पर, और भूमध्य रेखा पर अनंत लंबा। भूमध्य रेखा के निकट दाब-क्षेत्र बदलता है, और घर्षण कार्य करता है, इससे बहुत पहले कि पवन समदाब-रेखाओं के अनुदिश मुड़ सके: संतुलन कभी स्थापित नहीं होता, और वायु आंशिक रूप से प्रवणता के नीचे, मध्यम गति से बहती है।
\end{solution}

\begin{solution}{exo:b1:atmospheric-circulation:6}
(क)~$u = 464.6\,\sin^2\varphi/\cos\varphi$: \qty{14.2}{m/s} \ang{10} पर, \qty{32.2}{m/s} \ang{15} पर, \qty{57.8}{m/s} \ang{20} पर, \qty{134.1}{m/s} \ang{30} पर। (ख)~$u = 464.6\,(\cos^2\ang{10} - \cos^2\varphi)/\cos\varphi$: \qty{17.7}{m/s} \ang{15} पर, \qty{42.9}{m/s} \ang{20} पर, \qty{117.9}{m/s} \ang{30} पर। (ग)~भूमध्य रेखा से दूर शुरू होने वाली वायु अक्ष के निकट शुरू होती है, बचाने को कम कोणीय संवेग के साथ: वह हर \omterm{def:b1:spherical-astronomy:geographic}{अक्षांश} पर दुर्बल पछुआ पवन के साथ पहुँचती है। ऊपर की शाखा को फैलाने से अनुमानित जेट नीचे आता है, यद्यपि प्रेक्षित \qty{40}{m/s} तक पहुँचने के लिए पर्याप्त नहीं।
\end{solution}

\begin{solution}{exo:b1:atmospheric-circulation:7}
(क)~पवन उत्तर-पूर्व की ओर बहती है। उसकी ओर पीठ करके, उत्तर-पूर्व की ओर मुँह करके, निम्न बाईं ओर है: उत्तर-पश्चिम में। (ख)~पवन दक्षिण-पूर्व की ओर बहती है; दक्षिणी गोलार्ध में, पवन की ओर पीठ किए व्यक्ति की दाईं ओर निम्न होता है: दक्षिण-पश्चिम में। (ग)~सतह की पवन समदाब-रेखाओं से निम्न की ओर मुड़ती है, यानी उत्तरी गोलार्ध में \omterm{def:b1:atmospheric-circulation:geostrophic}{भू-विक्षेपी पवन} से वामावर्त। \omterm{def:b1:atmospheric-circulation:geostrophic}{भू-विक्षेपी पवन}, और उसकी बाईं ओर समकोण पर निम्न, इसलिए (क) के उत्तर से दक्षिणावर्त मुड़े हुए हैं: निम्न उत्तर-पश्चिम और उत्तर के बीच है, दिगंश $\ang{315} + \alpha$ पर, जहाँ $\alpha$ पार करने का कोण है।
\end{solution}

\begin{solution}{exo:b1:atmospheric-circulation:8}
(क)~भूमध्य रेखा और \ang{30} उ के बीच का क्षेत्रफल $2\pi a_\oplus^2\sin\ang{30} =
\qty{1.275e14}{m^2}$ है; दो स्तरों के बीच एक वर्ग मीटर के ऊपर वायु का द्रव्यमान $\Delta p/g = \num{9e4}/9.81 = \qty{9174}{kg}$ है; कुल
\qty{1.170e18}{kg}। (ख)~$\num{1.170e18}/\num{8.7e10} = \qty{1.34e7}{s}$, लगभग
\num{156} दिन, पाँच महीने। (ग)~$\num{1.170e18}/\num{21.4e10} = \qty{5.47e6}{s}$, लगभग \num{63} दिन, दो महीने।
\end{solution}

\begin{solution}{exo:b1:atmospheric-circulation:9}
(क)~$\num{0.01} \times \num{86400} = \qty{864}{m}$ एक दिन में। (ख)~$0.864 \times 9.8 =
\qty{8.5}{K}$ एक दिन में। (ग)~वायु अवरक्त विकिरण उत्सर्जित करके ठंडी होती है: वायुमंडल विकिरण से जितनी ऊर्जा खोता है उससे अधिक सोखता है, और वायु का नीचे उतरना, संपीड़न से उसे गरम करते हुए, उस हानि की भरपाई करता है। (घ)~नीचे उतरती वायु संपीड़न से गरम और सापेक्ष संतृप्ति के सापेक्ष शुष्क होती है; उसके मेघ वाष्पित हो जाते हैं और वर्षा नहीं बन सकती: \omterm{def:b1:atmospheric-circulation:cells}{उपोष्णकटिबंधीय उच्च} साफ़ आकाश और कम वर्षा की पेटियाँ हैं, जहाँ भूमि पर गरम मरुस्थल बनते हैं।
\end{solution}

\begin{solution}{exo:b1:atmospheric-circulation:10}
(क)~$V^2 + frV - fr|\vec v_g| = 0$ जहाँ $fr = \qty{55.85}{m/s}$:
$V = \tfrac12\bigl(-55.85 + \sqrt{55.85^2 + 4 \times 55.85 \times 15}\bigr) =
\qty{12.3}{m/s}$, \omterm{def:b1:atmospheric-circulation:geostrophic}{भू-विक्षेपी पवन} का \qty{82}{\percent}। (ख)~$V^2 - frV +
fr|\vec v_g| = 0$ के वास्तविक मूल तभी हैं यदि $(fr)^2 \geq 4fr|\vec v_g|$ है, यानी
$|\vec v_g| \leq fr/4 = \qty{14.0}{m/s}$। (ग)~\omterm{def:b1:atmospheric-circulation:cyclone}{प्रतिचक्रवात} के केंद्र के निकट, जहाँ $r$ छोटा है, सीमा $fr/4$ छोटी है: दाब-प्रवणता दुर्बल होनी चाहिए, समदाब-रेखाएँ दूर-दूर और पवनें हल्की। अवदाब की ऐसी कोई सीमा नहीं: उसकी समदाब-रेखाएँ केंद्र के निकट सघन हो सकती हैं, प्रबल पवनों के साथ।
\end{solution}

\begin{solution}{exo:b1:atmospheric-circulation:11}
(क)~छोटे $\varphi$ के लिए, $\sin\varphi \approx \varphi$ और $\cos\varphi \approx 1$, इसलिए $u \approx \Omega a_\oplus\varphi^2$। (ख)~$\ang{10} = \qty{0.1745}{rad}$:
$464.6 \times 0.1745^2 = \qty{14.15}{m/s}$, जबकि सटीक \qty{14.23}{m/s},
\qty{0.5}{\percent} कम। (ग)~$464.6\sin^2\varphi/\cos\varphi = 27.0$ का संख्यात्मक हल $\varphi = \ang{13.7}$ देता है। (घ)~वलय पहले ही \ang{14} से कम पर वास्तविक वार्षिक-माध्य जेट से मेल खाता है: कोणीय-संवेग का संरक्षण \omterm{def:b1:atmospheric-circulation:cells}{हैडली कोष्ठ} की ऊपरी शाखा का वर्णन केवल गहरे उष्णकटिबंध में करता है; लगभग \ang{15} के ध्रुव की ओर भँवर उसका स्थान ले लेते हैं।
\end{solution}

\begin{solution}{exo:b1:atmospheric-circulation:12}
(क)~माध्य $31.7/12 = \ang{2.6}$ उ; परिसर \ang{12.8} द से \ang{10.9} उ तक, यानी \ang{23.7}। (ख)~सबसे उत्तरी स्थिति सितंबर में आती है, जून संक्रांति से तीन महीने बाद; सबसे दक्षिणी फ़रवरी में, दिसंबर संक्रांति से दो महीने बाद। विलंब ऊष्मा-संचयन से आता है, सबसे अधिक महासागर की \omterm{def:b1:fluid-envelopes-structure:ocean-layers}{मिश्रित परत} में (\cref{ch:b1:fluid-envelopes-structure}), जो धीरे गरम और ठंडी होती है। (ग)~ज़ोनल माध्य सभी देशांतरों पर औसत लेता है: यदि ज़मीन के ऊपर वर्षा की पट्टी ग्रीष्मकालीन गोलार्ध में दूर तक झूलती है, तो सभी देशांतरों का माध्य भूमध्य रेखा को पार कर सकता है, भले ही किसी एक महासागर पर पट्टी ऐसा न करे।
\end{solution}

\begin{solution}{pb:b1:atmospheric-circulation:1}
\begin{enumerate}
  \item $\rho = p/(R_dT) = \num{100600}/(287.05 \times 283.15) =
        \qty{1.238}{kg.m^{-3}}$।
  \item $|\nabla p| = \qty{400}{Pa}/\qty{2.5e5}{m} = \qty{1.6e-3}{Pa.m^{-1}}$।
  \item $|\nabla p|/\rho = \num{1.6e-3}/1.238 = \qty{1.293e-3}{m.s^{-2}}$।
  \item $\Omega = 2\pi/\num{86164.2} = \qty{7.292e-5}{rad.s^{-1}}$;
        $f = 2\Omega\sin\ang{50} = \qty{1.117e-4}{s^{-1}}$।
  \item $|\vec v_g| = \num{1.293e-3}/\num{1.117e-4} = \qty{11.57}{m/s}$।
  \item $11.57/0.514444 = \qty{22.5}{kn}$; $11.57 \times 3.6 = \qty{41.7}{km/h}$।
  \item समय $|\vec v_g|/(|\nabla p|/\rho)$ है, और $|\vec v_g| = (|\nabla
        p|/\rho)/f$ है, इसलिए वह $1/f = \qty{8951}{s}$, \qty{2.49}{h} है; और $2\pi/f =
        \qty{15.6}{h}$। लगभग एक दिन या अधिक चलने वाले प्रतिरूप इसलिए संतुलन के निकट हैं।
  \item \omterm{def:b1:atmospheric-circulation:pgf}{दाब-प्रवणता बल} उत्तर की ओर, निम्न की ओर इंगित करता है। \omterm{def:b1:atmospheric-circulation:geostrophic}{भू-विक्षेपी पवन} समदाब-रेखाओं के अनुदिश बहती है, निम्न उसकी बाईं ओर: पूर्व की ओर, पश्चिम से, दिगंश \ang{270}: एक पछुआ पवन।
  \item पश्चिमी पवन की ओर पीठ करके, पूर्व की ओर मुँह किए, बाईं ओर उत्तर है, जहाँ निम्न है: नियम संतुष्ट है।
  \item स्टाफ़ स्टेशन-वृत्त से पश्चिम की ओर इंगित करता है, उस ओर जहाँ से पवन आती है; \qty{22.5}{kn} का \qty{20}{kn} तक पूर्णांकन होता है, दो पूर्ण पंख; वे स्टाफ़ के दक्षिणावर्त पक्ष पर खड़े हैं, उसकी उत्तरी ओर, निम्न की ओर।
  \item \ang{50} द पर, $f < 0$: पवन निम्न को अपनी दाईं ओर रखती है, इसलिए वह पश्चिम की ओर बहती है, पूर्व से।
  \item $f(\ang{30}) = \qty{7.292e-5}{s^{-1}}$: \qty{17.7}{m/s}; $f(\ang{60}) =
        \qty{1.263e-4}{s^{-1}}$: \qty{10.2}{m/s}।
  \item $f(\ang{5}) = \qty{1.271e-5}{s^{-1}}$ \qty{101.7}{m/s} देता है, ऐसी प्रवणताओं के साथ एक पवन जो कभी प्रेक्षित नहीं हुई। मोड़ने का समय $1/f$ \qty{22}{h} है: संतुलन स्थापित होने का समय नहीं है, और सूत्र लागू नहीं होता।
  \item $9.0/11.57 = 0.778$; समदाब-रेखाओं के अनुदिश $9.0\cos\ang{25} = \qty{8.16}{m/s}$, उनके आर-पार निम्न की ओर $9.0\sin\ang{25} = \qty{3.80}{m/s}$।
  \item \omterm{def:b1:atmospheric-circulation:geostrophic}{भू-विक्षेपी पवन} \ang{90} दिगंश की ओर बहती है; उत्तर में निम्न की ओर \ang{25} मुड़ी हुई, सतही पवन \ang{65} की ओर बहती है, \ang{245} से, पश्चिम-दक्षिण-पश्चिम।
  \item $\tan\alpha = 1/(f\tau)$ से $\tau = 1/(\num{1.117e-4}\tan\ang{25}) =
        \qty{1.92e4}{s}$, \qty{5.3}{h}; मॉडल $11.57\cos\ang{25} =
        \qty{10.5}{m/s}$ का पूर्वानुमान करता है, प्रेक्षित \qty{9.0}{m/s} से अधिक: वास्तविक घर्षण पवन में ठीक रैखिक नहीं है, और इस कोण पर उसे मॉडल से अधिक धीमा करता है।
  \item त्रिज्या $R = \qty{500}{km}$ और ऊँचाई $h =
        \qty{1}{km}$ के सिलिंडर से: $\dot m = 2\pi R h\rho u_r = 2\pi \times \num{5e5} \times 1000
        \times 1.238 \times 3.80 = \qty{1.48e10}{kg/s}$।
  \item $w = \dot m/(\rho\pi R^2) = 2hu_r/R = 2 \times 1000 \times 3.80/\num{5e5} =
        \qty{0.0152}{m/s}$, लगभग \qty{1310}{m} प्रति दिन।
  \item $1.31 \times 9.8 = \qty{12.9}{K}$ एक दिन में: वायु जल्द ही संतृप्ति तक पहुँचती है, और अवदाब मेघ और वर्षा से भर जाता है। \omterm{def:b1:atmospheric-circulation:cyclone}{प्रतिचक्रवात} के नीचे सतह की वायु बाहर की ओर बहती है, ऊपर की वायु नीचे उतरती और गरम होती है, और आकाश साफ़ होता है।
  \item $\Omega a_\oplus = \num{7.292e-5} \times \num{6.371e6} = \qty{464.6}{m/s}$
        भूमध्य रेखा पर; $464.6\cos\ang{30} = \qty{402.3}{m/s}$ \ang{30} उ पर।
  \item $u = 464.6\,\sin^2\varphi/\cos\varphi$: \qty{57.8}{m/s} \ang{20} उ पर और
        \qty{134.1}{m/s} \ang{30} उ पर।
  \item $u = 464.6\,(\cos^2\ang{10} - \cos^2\ang{30})/\cos\ang{30} =
        \qty{117.9}{m/s}$।
  \item $134.1/42.3 = 3.17$। मध्य-अक्षांशों के भँवर, अवदाब और \omterm{def:b1:atmospheric-circulation:cyclone}{प्रतिचक्रवात}, ऊपरी शाखा से कोणीय संवेग ले जाते हैं और उसे कहीं और ले जाते हैं; और कोष्ठ की वायु भूमध्य रेखा पर विराम में रहने के बजाय अक्षांशों की एक पट्टी के ऊपर उठती है। ज़ोनल माध्य भी देशांतरों के पार एक टेढ़े-मेढ़े जेट का औसत लेता है।
  \item $M/M_0 = (402.3 + 42.3)\cos\ang{30}/464.6 = 0.829$: वायु ने उस कोणीय संवेग का \qty{17}{\percent} खो दिया है जो वह एक मुक्त वलय के रूप में रखती।
  \item P पर \omterm{def:b1:atmospheric-circulation:geostrophic}{भू-विक्षेपी पवन} \qty{11.6}{m/s} है, पश्चिम से; ज़मीन पर मापी गई पवन \qty{9.0}{m/s} है, उसका लगभग चार-पाँचवाँ भाग, और घर्षण द्वारा निम्न की ओर \ang{25} मुड़ी हुई।
\end{enumerate}
\end{solution}
